\section*{Erd\H{o}s Problem 800: a linear Ramsey bound for subdivided graphs}

\paragraph{Statement and claim.}
If an \(n\)-vertex graph \(G\) has no adjacent pair of vertices
both of degree at least three, then \(R(G)\leq12n\).
We reconstruct Alon's proof:
\url{https://www.cs.tau.ac.il/~nogaa/PDFS/rams.pdf},
\emph{Subdivided graphs have linear Ramsey numbers},
Proposition 1.3. See \url{https://www.erdosproblems.com/800}.
The external input is the following theorem of
Sidorenko and, independently, Goddard--Kleitman:
every graph on \(2e+1\) vertices with independence number
at most two contains every \(e\)-edge graph without
isolated vertices as a subgraph.
It is stated as Theorem 2.1 in the cited primary source.

\paragraph{Separating the high-degree vertices.}
The vertices of degree at least three are independent.
Extend them to an inclusion-maximal independent set \(W\).
Every vertex outside \(W\) has a neighbour in \(W\) and
total degree at most two. Consequently each component
of \(G-W\) is either an isolated vertex or one edge.
An isolated vertex has one or two neighbours in \(W\);
each endpoint of a component edge has exactly one neighbour
in \(W\). The two neighbours may coincide.

Form a simple auxiliary graph \(H\) on \(W\):
join different \(w,w'\) whenever a component of \(G-W\)
provides a path between them with one or two internal vertices.
Distinct components have disjoint internal vertices,
so \(e(H)\leq n-|W|\leq n\).
Several components may provide the same auxiliary edge;
we keep only one copy in \(H\).

\paragraph{Finding a compatible image of \(W\).}
Consider any red-blue colouring of \(K_{12n}\).
Call a vertex red if its red degree is at least \(6n\),
and blue otherwise; a blue vertex has blue degree at least
\(6n\), since its red degree is at most \(6n-1\).
After interchanging colours, choose \(6n\) red vertices
and call their set \(U\).
On \(U\), form a graph \(T\) joining pairs with at least
\(2n\) common red neighbours in the original complete graph.

There cannot be three mutually nonadjacent vertices
\(u_1,u_2,u_3\) in \(T\). Otherwise their red neighbourhoods
would have union of size at least
\[
 3(6n)-3(2n-1)>12n,
\]
by the elementary lower inclusion--exclusion bound.
Thus \(\alpha(T)\leq2\).
Apply the external theorem to the nonisolated part of \(H\),
which has at most \(n\) edges.
There is ample room in \(T\) for its \(2e(H)+1\) required
vertices, and afterwards for all isolated vertices of \(H\).
We obtain distinct images \(w''\in U\) for all \(w\in W\),
such that every auxiliary edge has endpoints with
at least \(2n\) common red neighbours.

\paragraph{Embedding each remaining component.}
Process the components of \(G-W\) one by one.
Fewer than \(n\) vertices have been used at each stage.
An isolated component vertex with one neighbour \(w\in W\)
can be placed at any unused red neighbour of \(w''\);
there are at least \(6n\) before removing used vertices.
With two neighbours in \(W\), use an unused common red
neighbour, from a set of size at least \(2n\).

For an edge component \(uv\), let its neighbours in \(W\)
be \(w,w'\).
If \(w\ne w'\), more than \(n\) unused common red neighbours
of \(w'',w'''\) remain.
If this set has a red edge, use its ends for \(u,v\).
If it has no red edge, it contains a blue \(K_n\),
and hence a blue copy of \(G\).
If \(w=w'\), use instead at least \(n\) unused red neighbours
of \(w''\), with the same red-edge-or-blue-clique conclusion.
The additional red adjacencies possibly created by this
procedure are harmless: Ramsey copies are subgraphs,
not induced subgraphs.

\paragraph{Conclusion.}
Either the procedure finishes a red copy of \(G\), or it
finds a blue \(K_n\) and therefore a blue copy of \(G\).
Thus every colouring of \(K_{12n}\) contains the required
monochromatic graph. Isolated vertices of \(G\) are already
included in \(W\), so no connectedness assumption was used.
This proves the full assertion relative to the stated
triangle-versus-graph Ramsey theorem.

\paragraph{Completion estimate.}
\textbf{COMPLETION: 100\%}
